\begin{dissertationintroduction}
    
\setcounter{chapter}{0}
\chapter*{Introduction}

\label{chap:introduction}

Insulin resistance (IR) is a condition of decreased sensitivity to normal circulating levels of insulin in target tissues. It is the root cause of many epidemic diseases and conditions, including type 2 diabetes (T2D), cardiovascular disease (CVD), and metabolic syndrome (MetS)\cite{James2021-vu,Adeva-Andany2019-df,Park2024-da}. In addition to its clinical heterogeneity, there are also many insulin responsive tissues, including the liver, adipose tissue, and skeletal muscle, and many cell types that facilitate crosstalk between the different organs. Changes in one tissue or cell type can perturb whole body IR\cite{James2021-vu}.  

One class of drugs directly reduces IR by increasing insulin sensitivity. Thiazolidinediones (TZDs) decrease IR by agonizing PPAR$\gamma$\cite{Yki-Jarvinen2004-bc}, a nuclear hormone receptor encoded by gene \textit{PPARG}. TZDs have demonstrated clinical efficacy in treating T2D, CVD, and metabolic dysfunction-associated steatotic liver disease (MASLD), but their use has been limited by serious side effects, such as weight gain and fluid retention, which are a direct consequence of PPAR$\gamma$ agonism through interaction at its ligand-binding domain (LBD)\cite{Kernan2016-as,Musso2017-xi,Arnold2019-kq}. Much pharmacological development has been focused on the development of selective PPAR$\gamma$ modulators that retain the benefits of TZDs without the attendant adverse effects, but successful compounds have not reached the clinic\cite{DePaoli2014-lp,Higgins2010-yv}. Thus, the need for alternative approaches to therapeutically activate PPAR$\gamma$ without PPAR-mediated side effects remains unmet.  

Intriguingly, TZDs lead to weight gain by increasing a specific depot of adipose tissue that is metabolically beneficial\cite{Zhang2024-sl}. Previously, adipose tissue was regarded as only an energy reservoir, and it was understudied compared to other organs \cite{Sebo2019-nh}. We now know it is a dynamic endocrine organ that plays a central role in many physiological processes, including insulin sensitivity and immune responses\cite{Sakers2022-wb}. It is found throughout the human body, with tissue and organ specific depots, but white adipose tissue is generally categorized as either subcutaneous fat, located beneath the skin, or visceral fat, which surrounds the internal organs. Subcutaneous adipose tissue (SAT) is considered metabolically beneficial, while visceral adipose tissue (VAT) is unfavorable\cite{Mathur2022-fx}. Increased VAT is associated with higher risk of both T2D and CVD, whereas SAT is neutral or even protective\cite{Agrawal2023-ee}. While there have been many investigations on differential gene expression and cytokine secretion between SAT and VAT\cite{Mathur2022-fx,Ahn2019-us,Passaro2017-zk,Gesta2006-so,Tchkonia2007-xa,Wajchenberg2002-pt,Lee2010-fl}, the molecular basis of SAT and VAT heterogeneity is largely unknown. Further study of the mechanistic differences between adipocytes in SAT and VAT could elucidate how they contribute to metabolic health in opposing manners.   

While there are lifestyle factors that contribute to metabolic health, genetics also play a role in determining risk of IR and its associated diseases\cite{James2021-vu}. Genome-wide association studies (GWAS) have identified many loci associated with MetS\cite{Park2024-da}, but most GWAS do not include rare variants, which have a larger effect size\cite{Wang2021-au}. Moreover, after identification, the causal genes, mechanisms of action, and tissues of effect need to be elucidated for GWAS loci. Functional genomics studies how genotype affects phenotype at the genome-wide scale through the analysis of genes, RNA transcripts, proteins, and other biological components\cite{Farrell2017-lq}. Genome editing is a common method applied in functional genomics to both isolate causal variants from GWAS loci and enable characterization of the causal genes. The recent development of single cell CRISPR methods, such as Perturb-seq and CROP-seq, produce higher resolution cell phenotypes, generating whole transcriptomes of edited cells\cite{Datlinger2017-yc,Dixit2016-yg}. Combined with the hypotheses generated from GWAS, large screens, and cell atlases, single cell CRISPR and other functional genomics methods can help determine mechanisms underlying IR.   

In this dissertation, I present the following three chapters in which I leverage functional genomics and human genetics to study the role of \textit{PPARG} in insulin resistance and establish the framework to compare matched subcutaneous and visceral adipocytes in vitro.  

In Chapter 1, I perform a genome editing screen across \textit{PPARG} to investigate the effects of edits on each coding exon of the gene. This screen was performed with CRISPR-Cas9 in a human monocyte cell line, THP-1, which is a suitable model for large scale genetic perturbation that phenocopies \textit{PPARG} related transcriptional responses in adipocytes\cite{Majithia2016-bm,Agostini2018-bx}. We discovered that edits in the 5’ end of the gene still generated a functional isoform of PPAR$\gamma$, indicative of an alternative translational start site. These results were corroborated with analysis of an analogous screen that introduced the edited PPARG transcripts exogenously, and the findings were characterized with gene and protein expression. A truncated isoform, referred to as PPAR$\gamma$ M135, can be translated from methionine encoded by codon 135, and this isoform has higher transactivation of canonical PPAR$\gamma$ target genes in response to TZDs.  

In Chapter 2, I extend the findings of Chapter 1 to study the insulin sensitizing capacity of the novel PPAR$\gamma$ isoform. We generated adipocytes that express PPAR$\gamma$ M135 and found significantly increased insulin sensitizing activity in cells with PPAR$\gamma$ M135 as opposed to the endogenous PPAR$\gamma$2 isoform. Additionally, human carriers of rare variants in PPARG were evaluated for the severity of metabolic syndrome and insulin resistance, as measured by the triglyceride to HDL (TG/HDL) ratio. Human carriers of variants with decreased function in the N-terminal portion of PPAR$\gamma$, upstream of p.M135, have nominally lower metabolic syndrome severity and insulin resistance than carriers of variants with decreased function downstream of p.M135. Altogether, the data suggests that PPAR$\gamma$ M135 can be metabolically beneficial in humans as a result of increased insulin sensitization.   

In Chapter 3, I present the analytical framework of how to reprogram preadipocytes into depot-specific adipocytes. I identify factors to reprogram cells into either subcutaneous or visceral adipocytes using established methods and evaluate them on independent transcriptomic datasets. I further detail how to evaluate reprogrammed cells using pseudobulks of single cell/nuclei RNA-sequencing.   

In summary, my dissertation applies functional genomics to investigate genes and tissues that contribute to insulin resistance. Through genetic screens, transcriptomics, and proteomic analysis, I uncovered an alternative translational start site in PPARG that generates PPAR$\gamma$ M135, an isoform with increased transactivation in response to TZDs. PPAR$\gamma$ M135 results in increased insulin sensitivity, as measured in an adipocyte cell line and supported by analysis of human carriers. I additionally propose detailed methods for reprogramming and evaluating depot-specific adipocytes to deepen our understanding of adipocytes and enable further research into their function. In total, my dissertation aims to extend research on genes and cell types critical to preventing insulin resistance.

\end{dissertationintroduction}

\bibliographystyle{unsrtnat}
\bibliography{bib_0.bib}
